\chapter{Đào sâu: Từ điển trên MongoDB --- mô hình document, tìm kiếm regex và những đánh đổi thẳng thắn}
\label{ch:dd-mongo}

Chương~8 lập luận cho từ điển là service document của dự án: một mục từ
được đọc như một khối, ghi như một khối, và lồng nhau tự nhiên. Lập luận
đó vững. Chương này thử nó theo cách một kỹ sư senior hoài nghi sẽ làm
--- đọc các truy vấn service thực sự chạy, hỏi mỗi truy vấn tốn gì ở
mười nghìn từ mỗi người dùng, và thừa nhận nơi một cơ sở dữ liệu quan
hệ cũng làm tốt không kém. Mục tiêu không phải đảo ngược quyết định mà
để bạn bảo vệ được cả hai phía của nó trong phỏng vấn, với code ngay
trước mặt.

Ba câu hỏi sắp xếp chương này:

\begin{enumerate}
  \item \textbf{Mô hình document mua được gì ở đây}, và Postgres với
  \code{jsonb} sẽ tốn gì thay vào đó?
  \item \textbf{Truy vấn tìm kiếm thực sự thực thi thế nào} --- phần nào
  chạm index, phần nào quét, và trang được dựng ở đâu?
  \item \textbf{Chuyện gì gãy khi từ điển lớn lên} --- và trong số đó,
  cái nào là lỗi của Mongo, cái nào là lỗi của code?
\end{enumerate}

\section{Vì sao MongoDB ở đây, và nó giải quyết gì}

Mở \code{WordDefinition}. Nó là toàn bộ schema của service:

\begin{lstlisting}[language=Java]
@Document(collection = "word_definitions")
public class WordDefinition {
    @Id private String id;
    private String word;
    private String createdByUserId;                       (*@\co{1}@*)
    private String meaning;
    private String usage;
    private String notes;
    private List<String> examples;
    private List<String> tags;                            (*@\co{2}@*)
    /** Image attachments stored inline as base64 data URLs. */
    private List<String> images;                          (*@\co{3}@*)
    private Map<String, Object> extra;                    (*@\co{4}@*)
    @CreatedDate private Instant createdAt;
    @LastModifiedDate private Instant updatedAt;
}
\end{lstlisting}

\begin{description}
  \item[\co{1}] Mọi truy vấn trong service đều giới hạn theo trường
  này. Nó là hàng rào tenant và, như sẽ thấy, là thứ duy nhất index
  giúp được.
  \item[\co{2}] Hai mảng nhúng trong document. Trong Postgres đây là
  hai bảng con hoặc hai cột \code{text[]}; ở đây chúng là một phần của
  mục từ và không tốn gì để lấy cùng.
  \item[\co{3}] Ảnh dưới dạng chuỗi base64 inline. Đây là ``byte trong
  cơ sở dữ liệu'' của Chương~35 khoác áo document, kèm một trần cứng:
  một document BSON không được vượt 16\,MB. Một tá ảnh điện thoại và
  từ đó không lưu được nữa.
  \item[\co{4}] Một map tự do. Lối thoát hiểm của mô hình document: các
  trường code chưa biết tới. Cũng là lý do ``không schema'' cần kỷ luật
  (Chương~8).
\end{description}

Mô hình mua được gì: một lượt đọc trả về mọi thứ trang cần, không join,
không N+1 qua các bảng; thêm trường là thay đổi code, không phải
migration; và hình dạng trên đường truyền là hình dạng trên đĩa, khiến
service gần như toàn là keo dán. Nó tốn gì: quan hệ giữa các mục từ
(\code{WordLink}) không nằm trong document và phải duyệt bằng tay; mọi
trường mảng không index được theo tiền tố cho tìm kiếm chuỗi con; và
không gì ngăn một document lớn không giới hạn.

\begin{decision}[MongoDB, và lập luận cho việc lẽ ra chọn Postgres]
Từ điển là một sổ tay theo người dùng với các trường lồng nhau, tùy
chọn, tiến hóa --- sân nhà của mô hình document --- và không chia sẻ gì
với các service quan hệ, nên cách ly không tốn gì. Đó là lý do tốt.
Phản biện thật thà: Postgres có \code{jsonb} cho phần lồng nhau,
\code{text[]} với index GIN cho tag, \code{pg\_trgm} cho đúng kiểu tìm
chuỗi con service này chạy, và CTE đệ quy cho đồ thị từ, tất cả trong
một cơ sở dữ liệu đội đã vận hành hai lần. Câu trả lời senior cho ``vì
sao Mongo?'' vì thế không phải ``vì document''; mà là ``vì dữ liệu có
hình dạng document \emph{và} chúng tôi chấp nhận một cơ sở dữ liệu thứ
hai để service triển khai độc lập và để học polyglot persistence --- và
đây là cái giá''. Điều gì sẽ lật: yêu cầu join từ với khóa học hay
người dùng, hoặc một đội vận hành không muốn chạy engine thứ hai.
\end{decision}

\section{Đọc đường tìm kiếm}

Endpoint ai cũng dùng là tìm kiếm. Nó chạy qua ba mẩu code. Truy vấn
trong repository:

\begin{lstlisting}[language=Java]
@Query("{ 'createdByUserId': ?0, '$or': [ "
        + "{ 'word': { '$regex': ?1, '$options': 'i' } }, "
        + "{ 'tags': { '$regex': ?1, '$options': 'i' } } ] }")
List<WordDefinition> searchByWordOrTag(String userId, String regex);
\end{lstlisting}

Index đang tồn tại, từ \code{MongoIndexConfig}:

\begin{lstlisting}[language=Java]
mongoTemplate.indexOps("word_definitions")
        .ensureIndex(new Index().on("createdByUserId", Sort.Direction.ASC));
\end{lstlisting}

Và phương thức service biến danh sách thành trang:

\begin{lstlisting}[language=Java]
List<WordDefinitionResponse> responses = wordDefinitionRepository
        .searchByWordOrTag(userId, escapeRegex(query.trim()))   (*@\co{1}@*)
        .stream()
        .map(w -> toResponse(w, userId))                        (*@\co{2}@*)
        .sorted(Comparator.comparing(WordDefinitionResponse::word,
                String.CASE_INSENSITIVE_ORDER))                 (*@\co{3}@*)
        .toList();
int start = (int) pageable.getOffset();
int end = Math.min(start + pageable.getPageSize(), responses.size());
List<WordDefinitionResponse> pageContent = start < responses.size()
        ? responses.subList(start, end) : List.of();            (*@\co{4}@*)
return new PageImpl<>(pageContent, pageable, responses.size());
\end{lstlisting}

\begin{description}
  \item[\co{1}] Chữ người dùng gõ được ghép vào một biểu thức chính
  quy. \code{escapeRegex} thêm gạch chéo ngược trước mọi ký tự đặc biệt
  trước, nên \code{a.*} khớp đúng ba ký tự đó và \code{(} không làm
  server sập vì lỗi parse. Bỏ qua bước này là có regex injection: nhẹ
  thì kết quả sai, nặng thì một pattern thảm họa ghim chết một lõi CPU.
  \item[\co{2}] \code{toResponse} chạy thêm hai truy vấn mỗi từ --- cha
  và con từ \code{word\_links} --- nên một lượt tìm khớp $n$ từ phát ra
  $2n + 1$ vòng đi-về. Đây là bài toán N+1, tái sinh trong một cơ sở dữ
  liệu vốn được cho là tránh nó.
  \item[\co{3}] Sắp xếp diễn ra trong JVM, trên \emph{mọi} kết quả khớp,
  không phải trên trang.
  \item[\co{4}] Phân trang cũng vậy. Mongo trả về mọi kết quả; service
  giữ hai mươi và vứt phần còn lại. Trang 50 tốn đúng bằng trang 1,
  cộng thêm việc sao chép.
\end{description}

\subsection{Server thực thi nó thế nào}

Hỏi chính Mongo. Với stack đang chạy:

\begin{handson}[đọc kế hoạch thực thi]
\begin{lstlisting}[language=bash]
$ docker compose exec mongodb mongosh ts_dictionary --quiet --eval '
  db.word_definitions.find({ createdByUserId: "1", $or: [
    { word: { $regex: "clou", $options: "i" } },
    { tags: { $regex: "clou", $options: "i" } } ] })
  .explain("executionStats").executionStats'
\end{lstlisting}
Lược bớt, câu trả lời trông thế này:
\begin{lstlisting}
{ nReturned: 3,
  totalKeysExamined: 412,
  totalDocsExamined: 412,
  executionStages: {
    stage: "FETCH",
    filter: { $or: [ { word: { $regex: "clou", $options: "i" } },
                     { tags: { $regex: "clou", $options: "i" } } ] },
    inputStage: {
      stage: "IXSCAN",
      indexName: "createdByUserId_1",
      indexBounds: { createdByUserId: [ "[\"1\", \"1\"]" ] } } } }
\end{lstlisting}
Ba dòng kể trọn câu chuyện: \code{IXSCAN} trên \code{createdByUserId}
thu hẹp về 412 document của người dùng này; \code{FETCH} rồi nạp từng
document từ đĩa; \code{filter} chạy cả hai regex trên mỗi document. Tỷ
số \code{totalDocsExamined / nReturned} là 137 --- mỗi từ trả về, 137
từ được đọc rồi loại.
\end{handson}

Vì sao regex không dùng được index: B-tree sắp xếp chuỗi theo các ký
tự đầu. \code{\^{}clou} (có neo, phân biệt hoa thường) ánh xạ thành một
khoảng key và Mongo sẽ dùng index trên \code{word} cho nó. \code{clou}
không neo có thể xuất hiện bất kỳ đâu trong chuỗi, và \code{i} nghĩa là
thứ tự byte của index thậm chí không khớp phép so sánh; không có khoảng
nào để nhảy tới. Bộ lập kế hoạch của Mongo nói rõ điều này: một
\code{\$regex} không phân biệt hoa thường hoặc không neo là quét toàn
bộ phần mà các điều kiện khác để lại. Ở đây đó là lát của người dùng,
và vì thế index \code{createdByUserId} đang làm toàn bộ việc hữu ích,
còn index kép đề xuất về sau không thay đổi phần regex --- nó thay đổi
điều xảy ra \emph{sau} đó.

\subsection{Điều kiện trên mảng}

\code{tags} là một mảng. Một điều kiện vô hướng trên trường mảng khớp
nếu \emph{bất kỳ} phần tử nào khớp --- không cần \code{\$elemMatch} cho
một điều kiện đơn. Đó là tiện lợi thật so với SQL, nơi cùng điều kiện
là một \code{EXISTS} vào bảng con. Cũng vì thế regex chạy một lần cho
mỗi tag mỗi document: một document có tám tag tốn tám lần đánh giá
pattern bên cạnh một lần cho \code{word}.

\section{Đồ thị}

\code{WordLink} là một danh sách kề: một document mỗi cạnh với
\code{parentWordId}, \code{childWordId}, \code{position}. Index ở cả
hai đầu là đúng --- ``con của X'' và ``cha của X'' mỗi cái là một lượt
tra index. Chi phí nằm ở việc duyệt:

\begin{lstlisting}[language=Java]
private WordGraphResponse buildGraphNode(String wordId, String userId,
                                         Set<String> visited) {
    if (visited.contains(wordId)) return null;
    visited.add(wordId);
    WordDefinition word = wordDefinitionRepository
            .findByIdAndCreatedByUserId(wordId, userId).orElse(null); (*@\co{1}@*)
    if (word == null) return null;
    List<WordLink> childLinks = wordLinkRepository
            .findByParentWordIdOrderByPosition(wordId);              (*@\co{2}@*)
    List<WordGraphResponse> children = childLinks.stream()
            .map(link -> buildGraphNode(link.getChildWordId(), userId,
                    visited))
            .filter(Objects::nonNull).toList();
    return new WordGraphResponse(word.getId(), word.getWord(),
            word.getMeaning(), children);
}
\end{lstlisting}

\begin{description}
  \item[\co{1}] Một truy vấn mỗi node.
  \item[\co{2}] Thêm một nữa mỗi node. Một đồ thị 2{,}000 từ tốn
  4{,}000 vòng đi-về, tuần tự, trên thread request. Mỗi lượt dưới một
  mili giây trên cùng host; qua một chặng mạng là một trang nhiều giây.
  Nhúng \code{childIds} vào document cha sẽ biến nó thành hai truy vấn
  tổng cộng (\code{find} mọi từ của người dùng, dựng trong bộ nhớ) ---
  với cái giá cập nhật hai document mỗi liên kết, mà không có
  transaction giữ chúng cùng nhau.
\end{description}

Mongo có một toán tử đồ thị, \code{\$graphLookup}, duyệt
\code{word\_links} đệ quy phía server trong một aggregation. Nó trả về
tập phẳng các node với tới được kèm trường độ sâu, và cây được dựng lại
trong bộ nhớ. Tập \code{visited} của dự án làm bằng tay điều
\code{\$graphLookup} làm với \code{maxDepth}; khác biệt là $2n$ vòng
đi-về so với một.

\section{Transaction và nhất quán, thật thà}

Mongo hỗ trợ transaction ACID nhiều document từ bản 4.0 --- trên một
replica set. Dự án chạy một \code{mongod} đơn (Compose lẫn
StatefulSet), và một node standalone không thể bắt đầu transaction nào.
Nên \code{deleteWord} là hai lượt ghi độc lập: xóa định nghĩa, rồi
\code{deleteByParentWordIdOrChildWordId}. Crash giữa hai lượt để lại
các cạnh trỏ vào một từ đã mất; việc duyệt chịu được
(\code{orElse(null)}, rồi lọc), là tư thế phòng thủ đúng, nhưng là một
hợp đồng toàn vẹn dữ liệu ẩn. Cách sửa không tốn dòng code nào và một
dòng vận hành là chạy replica set một node (\code{--replSet rs0} cộng
một lời gọi initiate): transaction trở nên khả dụng, và cả change
stream bạn sẽ muốn cho mẫu outbox của Chương~11.

Write concern và read preference quan trọng ngay khi có node thứ hai.
Mặc định \code{w: majority} nghĩa là một lượt ghi được xác nhận đã nằm
trên đa số node; đọc từ secondary có thể trả về trạng thái cũ hơn. Với
một sổ tay một người dùng, đọc-được-điều-mình-vừa-ghi trên primary là
điều người dùng mong đợi, và là mặc định --- nói vậy trong phỏng vấn,
và rằng bạn sẽ xem lại nó trước khi thêm read replica.

\section{Ưu và nhược, thẳng thắn}

\begin{center}
\begin{tabular}{@{}p{0.3\linewidth}p{0.62\linewidth}@{}}
\toprule
\textbf{Điểm mạnh} & \textbf{Dự án trả giá ở đâu} \\
\midrule
Một document, một lượt đọc, không join &
  Liên kết nằm riêng; đồ thị là $2n$ truy vấn \\
Mảng là hạng nhất (\code{tags}, \code{examples}) &
  Tìm chuỗi con trên chúng là quét từng phần tử \\
Không migration cho trường mới &
  \code{extra} và \code{images} lớn không kiểm soát tới 16\,MB \\
Tìm kiếm regex chạy ngay không cần gì &
  Nó không dùng được index trừ khi có neo và phân biệt hoa thường \\
Phân trang chỉ cách một \code{skip/limit} &
  Đường tìm kiếm phân trang trong Java sau khi nạp tất cả \\
Cơ sở dữ liệu độc lập, triển khai độc lập &
  Không auth trên \code{mongod} trong Compose \emph{lẫn} cụm \\
\bottomrule
\end{tabular}
\end{center}

Dòng cuối không phải đặc tính của Mongo mà là một sự thật của dự án
đáng biết trước khi ai đó hỏi: image chính thức chỉ bật xác thực khi
đặt thông tin root, và cả hai manifest đều không đặt. Trong một mạng
chỉ ClusterIP đó là rủi ro có cân nhắc, được ghi trong
\code{mongodb.yaml}; nó là thứ đầu tiên phải đổi nếu service có ngày
chia cụm với bất kỳ thứ gì không đáng tin.

\section{Chuyện gì gãy}

\begin{pitfall}[tìm kiếm chậm dần theo mỗi từ, rồi hết giờ]
Triệu chứng: độ trễ tìm kiếm tăng tuyến tính theo từ điển của người
dùng; tới vài nghìn từ, timeout của frontend nổ và trang không hiện gì.
Nguyên nhân: ba hệ số nhân chồng lên cùng một request --- regex quét
mọi document trong lát của người dùng; \code{toResponse} rồi phát hai
truy vấn mỗi kết quả; và sắp xếp cùng phân trang diễn ra sau cả hai. Cơ
sở dữ liệu đang làm đúng điều được yêu cầu. Sửa theo thứ tự hiệu quả:
phân trang trong cơ sở dữ liệu (bên dưới), gộp các lượt tra liên kết
thành một truy vấn \code{\$in} mỗi trang, và neo hoặc thay thế regex.
\end{pitfall}

\begin{pitfall}[document quá lớn]
Triệu chứng: lưu một từ thất bại với
\code{BsonMaximumSizeExceededException: ... exceeds 16793600 bytes}.
Nguyên nhân: \code{images} giữ các data URL base64, và base64 phình
thêm một phần ba; năm ảnh 3\,MB là document vượt giới hạn. Sửa: lưu ảnh
trong MinIO dưới tiền tố \code{dictionary/\{userId\}/\{wordId\}/} và
giữ key trong document --- quy tắc của Chương~35, áp dụng ở đây. Trong
lúc chờ, giới hạn mảng bằng validation ở
\code{WordDefinitionRequest}.
\end{pitfall}

\begin{pitfall}[regex không bao giờ trả về]
Triệu chứng: một request ghim một lõi ở 100\,\% và các truy vấn còn lại
của service xếp hàng sau nó. Nguyên nhân: một pattern như
\code{(a+)+\$} trên chuỗi dài --- backtracking thảm họa.
\code{escapeRegex} ngăn người dùng gửi một cái như vậy; rủi ro quay lại
vào ngày ai đó thêm ``tìm kiếm nâng cao'' truyền pattern thô, hoặc quên
escape ở một endpoint mới. Phòng thủ nhiều lớp: một \code{maxTimeMS}
trên truy vấn (\code{@Meta(maxExecutionTimeMs)} của Spring Data) để
pattern chạy hoang bị giết phía server sau một giây.
\end{pitfall}

\section{Chuyện gì gãy lúc 3 giờ sáng}

Node của StatefulSet MongoDB khởi động lại; pod ở trạng thái pending
bốn phút trong khi claim gắn lại.

\begin{itemize}
  \item \textbf{dictionary-service} trả về 500 trên mọi endpoint.
  Readiness probe của nó (Chương~17) không có kiểm tra cơ sở dữ liệu,
  nên gateway vẫn định tuyến tới nó. Người dùng thấy lỗi thay vì một
  trang ``dịch vụ không sẵn sàng''.
  \item \textbf{Mọi thứ khác} không bị chạm. Không service nào khác giữ
  kết nối Mongo, không service nào gọi dictionary-service, và không
  topic Kafka nào phụ thuộc vào nó. Đây là mỗi-service-một-cơ-sở-dữ-liệu
  hoạt động đúng thiết kế: bán kính thiệt hại của sự cố là đúng một
  tính năng.
  \item \textbf{Khi pod trở lại}, connection pool của driver tự kết nối
  lại; không cần khởi động lại và không có cache phải làm ấm.
  Create-on-write của Mongo nghĩa là một claim \emph{sai} sẽ lên với dữ
  liệu rỗng và im lặng (cạm bẫy của Chương~8) --- kiểm tra
  \code{show dbs} trước khi tin một báo cáo ``mất dữ liệu''.
\end{itemize}

Cái 3 giờ sáng còn lại thì chậm hơn: claim 5\,GB đầy. Mongo không từ
chối gọn một lượt ghi ở thời điểm đó; nó khựng lại, journal thất bại,
và process có thể thoát. Vì document có thể lớn lên (ảnh, \code{extra}),
không có trần nào trong schema ngăn được. Cảnh báo thuộc về mức dùng của
PVC, không phải ứng dụng.

\section{Nâng cấp giải pháp}

Bốn thay đổi, mỗi cái trọn vẹn, theo thứ tự đem lại hiệu quả.

\subsection{Phân trang trong cơ sở dữ liệu}

Thay truy vấn trả về danh sách bằng truy vấn trả về trang. Spring Data
Mongo cần một truy vấn đếm riêng cho trang từ \code{@Query}:

\begin{lstlisting}[language=Java]
@Query(value = "{ 'createdByUserId': ?0, '$or': [ "
        + "{ 'word': { '$regex': ?1, '$options': 'i' } }, "
        + "{ 'tags': { '$regex': ?1, '$options': 'i' } } ] }",
       count = true)
long countByWordOrTag(String userId, String regex);

@Query("{ 'createdByUserId': ?0, '$or': [ "
        + "{ 'word': { '$regex': ?1, '$options': 'i' } }, "
        + "{ 'tags': { '$regex': ?1, '$options': 'i' } } ] }")
Page<WordDefinition> searchByWordOrTag(String userId, String regex,
                                       Pageable pageable);          (*@\co{1}@*)
\end{lstlisting}

\begin{lstlisting}[language=Java]
public Page<WordDefinitionResponse> searchWords(String userId,
        String query, Pageable pageable) {
    Pageable sorted = PageRequest.of(pageable.getPageNumber(),
            pageable.getPageSize(),
            Sort.by("word").ascending());                            (*@\co{2}@*)
    Page<WordDefinition> page = (query == null || query.isBlank())
            ? wordDefinitionRepository.findByCreatedByUserId(userId, sorted)
            : wordDefinitionRepository.searchByWordOrTag(userId,
                    escapeRegex(query.trim()), sorted);
    return toResponses(page, userId);                                (*@\co{3}@*)
}
\end{lstlisting}

\begin{description}
  \item[\co{1}] Spring Data nối \code{skip}, \code{limit} và
  \code{sort} vào truy vấn. Regex vẫn quét lát của người dùng, nhưng chỉ
  hai mươi document được nạp, tuần tự hóa và ánh xạ.
  \item[\co{2}] Sắp xếp chuyển sang server. Thứ tự không phân biệt hoa
  thường cần một collation trên truy vấn hoặc collection
  (\code{\{locale: "en", strength: 2\}}); thêm nó vào index bên dưới để
  sort dùng được.
  \item[\co{3}] Các lượt tra liên kết thành một lô mỗi trang, tiếp theo.
\end{description}

\subsection{Gộp lô các liên kết}

\begin{lstlisting}[language=Java]
private Page<WordDefinitionResponse> toResponses(Page<WordDefinition> page,
                                                 String userId) {
    List<String> ids = page.map(WordDefinition::getId).toList();
    Map<String, List<String>> parents = new HashMap<>();
    Map<String, List<String>> children = new HashMap<>();
    for (WordLink l : wordLinkRepository.findByChildWordIdIn(ids)) {   (*@\co{1}@*)
        parents.computeIfAbsent(l.getChildWordId(), k -> new ArrayList<>())
               .add(l.getParentWordId());
    }
    for (WordLink l : wordLinkRepository
            .findByParentWordIdInOrderByPosition(ids)) {
        children.computeIfAbsent(l.getParentWordId(), k -> new ArrayList<>())
                .add(l.getChildWordId());
    }
    return page.map(w -> toResponse(w,
            parents.getOrDefault(w.getId(), List.of()),
            children.getOrDefault(w.getId(), List.of())));
}
\end{lstlisting}

\begin{description}
  \item[\co{1}] Hai truy vấn \code{\$in} mỗi trang thay vì hai mỗi từ:
  $2n + 1$ thành $3$. Cả hai dùng index sẵn có trên \code{childWordId}
  và \code{parentWordId}.
\end{description}

\subsection{Index cho sort, và tìm kiếm có neo}

\begin{lstlisting}[language=Java]
// MongoIndexConfig
mongoTemplate.indexOps("word_definitions").ensureIndex(
        new CompoundIndexDefinition(
                new Document("createdByUserId", 1).append("word", 1))
        .collation(Collation.of("en").strength(2)));               (*@\co{1}@*)
\end{lstlisting}

\begin{description}
  \item[\co{1}] Với index này, ``từ của người dùng này theo thứ tự chữ
  cái'' là một lượt đi dọc index, nên danh sách không lọc và mọi trang
  đã sort tránh được sort trong bộ nhớ. Và tìm theo \emph{tiền tố} ---
  \code{\{word: \{\$regex: "\^{}clou"\}\}} với cùng collation --- thành
  một lượt seek khoảng trên cùng index. Với từ điển, tiền tố thường là
  điều người dùng muốn khi gõ.
\end{description}

\subsection{Text index cho tìm-từ-ở-bất-kỳ-đâu}

Khi thực sự cần chuỗi con ở bất kỳ đâu, thôi dùng regex:

\begin{lstlisting}[language=Java]
// MongoIndexConfig: one text index per collection, weighted
mongoTemplate.indexOps("word_definitions").ensureIndex(
        new TextIndexDefinition.TextIndexDefinitionBuilder()
                .onField("word", 10F)                                (*@\co{1}@*)
                .onField("tags", 5F)
                .onField("meaning", 1F)
                .build());
\end{lstlisting}

\begin{lstlisting}[language=Java]
// DictionaryService
public Page<WordDefinitionResponse> textSearch(String userId,
        String terms, Pageable pageable) {
    TextCriteria text = TextCriteria.forDefaultLanguage()
            .matchingAny(terms.split("\\s+"));                       (*@\co{2}@*)
    Query q = TextQuery.queryText(text).sortByScore()
            .addCriteria(Criteria.where("createdByUserId").is(userId))
            .with(pageable);
    List<WordDefinition> hits = mongoTemplate.find(q, WordDefinition.class);
    long total = mongoTemplate.count(Query.of(q).skip(0).limit(0),
            WordDefinition.class);
    return toResponses(new PageImpl<>(hits, pageable, total), userId);
}
\end{lstlisting}

\begin{description}
  \item[\co{1}] Text index tách token và rút gốc từ; \code{cloud} khớp
  \code{clouds} và \code{clouded}, và một kết quả trong \code{word} xếp
  trên cùng kết quả đó trong \code{meaning}. Nó không khớp \emph{bên
  trong} một token (\code{lou} không tìm ra \code{cloud}) --- đó là
  đánh đổi.
  \item[\co{2}] Không regex, không escape, không backtracking. Index
  được hỏi trước; \code{createdByUserId} lọc các kết quả.
\end{description}

\begin{decision}[tiền tố có neo so với full text]
Regex neo tiền tố trên index kép trả lời ``từ bắt đầu bằng\ldots''
trong vài micro giây và không cần loại index mới; nó là mặc định đúng
cho ô gõ-là-gợi-ý. Text index trả lời ``mục từ về\ldots'' với xếp hạng
và rút gốc, với cái giá là index lớn hơn và không khớp giữa từ. Regex
không neo không phân biệt hoa thường --- thứ dự án có --- trả lời ``bất
kỳ thứ gì chứa\ldots'' và trả một lượt quét cho nó. Khuyến nghị là tiền
tố cho ô tìm kiếm, text search sau một nút ``tìm mọi nơi'' riêng, và
regex về hưu. Atlas Search (Lucene dưới API của Mongo) cho cả hai cùng
lúc, nhưng chỉ trên Atlas; trên một \code{mongod} tự host, hai công cụ
này là tất cả.
\end{decision}

\section{Ngoài dự án}

Mọi cơ sở dữ liệu document đều có hình dạng của chương này: DynamoDB
với mẫu truy cập khóa-phân-vùng-trước và hoàn toàn không có tìm chuỗi
con; Couchbase với N1QL và một index full-text; Elasticsearch là
sidecar mà các kho document mọc ra khi tìm kiếm trở thành sản phẩm.
\code{jsonb} cộng \code{pg\_trgm} của Postgres là câu trả lời phần lớn
đội nên cân nhắc trước bất kỳ thứ nào trong số đó --- nó biến đúng
\code{ILIKE '\%clou\%'} của service này thành một lượt seek index GIN.
Cách đặt vấn đề senior cho ``SQL hay NoSQL?'' không bao giờ là engine;
mà là các mẫu truy cập viết ra trước, rồi engine nào có index khớp với
chúng.

\section{Câu hỏi từ phỏng vấn thật}

\subsection*{Q: Tìm kiếm từ điển chậm. Bạn kiểm tra gì trước?}

Kế hoạch thực thi, trước mọi thứ: \code{explain("executionStats")}
trên đúng truy vấn service gửi, và hai con số từ nó ---
\code{totalDocsExamined} so với \code{nReturned}, và stage nào mang
\code{filter}. Ở service này nó cho thấy một \code{IXSCAN} trên
\code{createdByUserId} rồi một \code{FETCH} xem xét hàng trăm document
mỗi kết quả, vì regex không neo không phân biệt hoa thường không dùng
được index nào. Rồi nhìn điều xảy ra \emph{sau} truy vấn: ở đây, hai
lượt tra liên kết mỗi kết quả và sort cùng phân trang trong bộ nhớ.
Tôi sẽ sửa theo thứ tự đó --- phân trang trong cơ sở dữ liệu, gộp lô
liên kết bằng \code{\$in}, neo regex thành tiền tố trên index kép
\code{\{createdByUserId, word\}} với collation không phân biệt hoa
thường --- và đọc lại kế hoạch sau mỗi bước. Nếu ``bất kỳ đâu trong
từ'' là yêu cầu thật, đó là một text index, không phải regex nhanh hơn.

\subsection*{Q: Khi nào bạn không chọn MongoDB?}

Khi mẫu truy cập có tính quan hệ: dữ liệu đọc theo cách nhóm khác với
cách ghi, các bất biến trải qua nhiều thực thể (một bài nộp cần bài tập
và học sinh của nó tồn tại), hay báo cáo join qua cả miền. Khi đội đã
vận hành Postgres tốt và ``document'' chỉ cách một cột \code{jsonb}.
Khi transaction nhiều document là chuyện thường ngày chứ không phải
ngoại lệ --- Mongo có, nhưng trên replica set, với ngữ nghĩa retry bạn
phải tự code. Và khi dữ liệu sẽ được tìm theo chuỗi con ở quy mô lớn:
\code{pg\_trgm} là một dòng index; Mongo cần Atlas Search hoặc một
engine ngoài. Từ điển của dự án qua được các bài kiểm tra đó vì nó theo
người dùng, lồng nhau, và không join gì; miền khóa học trượt cả ba, và
đó là lý do nó ở Postgres.

\subsection*{Q: Phân trang một trăm nghìn kết quả mà không nạp chúng.}

Không bao giờ bằng \code{skip} ở độ sâu lớn --- \code{skip(90000)} vẫn
đi qua chín mươi nghìn mục index. Dùng phân trang \emph{keyset}: sort
trên một khóa có index và gần như duy nhất (\code{word}, phân định bằng
\code{\_id}) và hỏi \code{\{word: \{\$gt: lastWord\}\}} với
\code{limit(20)}; mỗi trang là một lượt seek index bất kể độ sâu. Trả
về khóa cuối làm cursor mờ thay vì số trang. Chỉ giữ truy vấn đếm nếu
UI thực sự cần tổng, và cache nó, vì đếm một regex có lọc là lại chính
lượt quét đó. Phân trang \code{subList} của dự án là ngược lại tất cả
những điều này: nó nạp mọi thứ và trả về hai mươi.

\subsection*{Q: Một service dùng Mongo, một service dùng Postgres. Giữ dữ liệu của chúng nhất quán thế nào?}

Bạn không làm được, theo nghĩa transaction --- không có two-phase
commit qua chúng và bạn cũng không nên muốn. Bạn giữ chúng nhất quán
\emph{cuối cùng} qua sự kiện (Chương~11): service sở hữu commit thay
đổi của mình và công bố một sự thật, service kia tiêu thụ và cập nhật
bản sao của riêng nó, và một bảng outbox ở phía sở hữu khiến ``đã
commit'' và ``đã công bố'' là một. Nơi dự án cần điều này hôm nay là id
người dùng: dictionary-service lưu \code{createdByUserId} là một chuỗi
do auth-service phát, và không bao giờ cần join trên nó, nên không cần
đồng bộ gì. Nếu có ngày cần tên hiển thị của người dùng, câu trả lời là
một consumer \code{UserRegistered} ghi vào một collection \code{users}
nhỏ --- không phải một lời gọi Feign mỗi từ, và không phải một cơ sở dữ
liệu dùng chung.

\subsection*{Q: Hình dạng document đã đổi. Migrate không gián đoạn.}

Mở rộng, di chuyển, thu hẹp. Triển khai code đọc được cả hai hình dạng
và ghi hình dạng mới (một alias \code{@Field} hay một converter xử lý
việc đổi tên); chạy một job nền viết lại document cũ theo lô với
\code{updateMany} trên một khoảng \code{\_id} có giới hạn, có tiết lưu,
lũy đẳng, tiếp tục được; rồi, khi đếm hình dạng cũ về không, triển khai
code chỉ đọc hình dạng mới. Mongo làm bước đầu dễ vì không gì cưỡng chế
schema, và nguy hiểm cũng vì lý do đó --- thêm một JSON Schema
validator trên collection ở cuối để hình dạng cũ không quay lại được.
Với trường \code{images} trong dự án này, migration còn là một lần
chuyển dữ liệu: viết lại mỗi data URL thành một object MinIO và thay
phần tử mảng bằng key của nó.

\subsection*{Q: Mô hình hóa đồ thị từ --- nhúng hay tham chiếu, và duyệt tốn gì?}

Tham chiếu, như dự án làm, khi cạnh có thuộc tính riêng
(\code{position}) và cả hai chiều đều được truy vấn; nhúng một mảng
\code{childIds} khi đồ thị nông, chủ yếu đọc, và cha-tới-con là chiều
duy nhất quan trọng. Chi phí duyệt là phần thật thà: với tham chiếu và
đệ quy ngây thơ đó là hai vòng đi-về mỗi node --- 4{,}000 cho đồ thị
2{,}000 từ --- chính là điều \code{buildGraphNode} làm. Cách sửa, theo
thứ tự, là nạp mọi cạnh của người dùng trong một truy vấn và duyệt
trong bộ nhớ; hoặc một \code{\$graphLookup} với \code{maxDepth}; hoặc
nhúng. Chu trình cần một tập visited dù đường nào; dự án có sẵn. Và nói
to đánh đổi: nhúng biến một lần tạo liên kết thành hai lượt ghi
document mà không có transaction trên node standalone.

\begin{quiz}
\begin{enumerate}
  \item Từ kết quả \code{explain}, \code{totalDocsExamined} là 412 và
  \code{nReturned} là 3. Nêu stage đang loại bỏ, nói vì sao index
  \code{createdByUserId} không giúp được nó, và nêu một thay đổi duy
  nhất trên pattern cho phép index trên \code{word} phục vụ nó.
  \item \code{searchWords} với \code{page=50\&size=20} trên người dùng
  có 3{,}000 từ: Mongo trả về bao nhiêu document, bao nhiêu truy vấn
  \code{word\_links} chạy, và trang được dựng ở đâu?
  \item \code{deleteWord} phát hai lượt ghi không có transaction. Giải
  thích vì sao transaction không khả dụng trên triển khai này, crash
  giữa hai lượt để lại gì, và dòng nào của \code{buildGraphNode} chịu
  được điều đó.
  \item Một người dùng đính sáu ảnh vào một từ và lưu thất bại. Nêu
  exception, phép tính tạo ra nó, và quy tắc lưu trữ từ Chương~35 sửa
  được nó.
  \item Với text index trên \code{word}, \code{tags} và
  \code{meaning}, truy vấn ``lou'' không trả về gì dù ``cloud'' tồn
  tại. Giải thích, và nói loại index nào trả lời được.
  \item MongoDB chết bốn phút lúc 3 giờ sáng. Liệt kê cái gì hỏng, cái
  gì không, và một thay đổi probe duy nhất biến 500 thành một ``chưa
  sẵn sàng'' sạch sẽ.
\end{enumerate}
\end{quiz}
